\documentclass[12pt,a4paper,oneside]{article}
\usepackage[brazil]{babel}
\usepackage[utf8]{inputenc}
\usepackage{olymp}
\usepackage{color}
\usepackage[dvipsnames]{xcolor}
\usepackage{listings}

\setcounter{problem}{1}

\begin{document}

\begin{problem}{Frequência}{freq}{2}
 
Faça um programa que leia uma matriz de dimensões $M \times N$ e um valor $V$ e informe quantas vezes este valor aparece em cada linha da matriz.

%\Notes
%
%Observações, se tiver.

\InputFile

A primeira linha contém dois
números inteiros $M$ e $N$, indicando respectivamente o número de linhas e de colunas da matriz.

As $M$ linhas seguintes contêm $N$ inteiros cada, indicando os
elementos de cada linha da matriz.

A última linha contém um único inteiro, o valor $V$.

Restrições: $1 \leq M,N \leq 1000$.

\OutputFile

Seu programa deve gerar $M$ linhas na saída, uma para cada linha da matriz, informando o número de vezes que o valor $V$ aparece naquela linha da matriz.


%\Notes
%Nesta questão, seu programa deve usar a função \texttt{fatorial} dada %acima exatamente como está, não a modifique! Sua
%tarefa é apenas criar o programa com a função \texttt{main}.

\Example

\begin{example}
\exmp{
4 5
1 2 3 2 1
2 3 4 3 2
3 4 3 2 3
0 1 2 1 0
3
}{
1
2
3
0
}%
\end{example}

%Comentários sobre o exemplo, se necessário.

\end{problem}

\end{document}
